\documentclass[11pt]{article}
\usepackage[margin=1in]{geometry}
\usepackage{amsmath,amssymb,amsthm,mathtools}
\usepackage{hyperref}
\usepackage{enumitem}

\newtheorem{theorem}{Theorem}
\newtheorem{lemma}[theorem]{Lemma}
\newtheorem{proposition}[theorem]{Proposition}
\newtheorem{corollary}[theorem]{Corollary}
\theoremstyle{definition}
\newtheorem{definition}[theorem]{Definition}
\newtheorem{remark}[theorem]{Remark}

\DeclareMathOperator{\ImOp}{Im}
\DeclareMathOperator{\ReOp}{Re}
\DeclareMathOperator{\supp}{supp}
\let\Im\ImOp
\let\Re\ReOp

\title{Hilbert Space Identification of the Zeta Spectral Process}
\author{Stephen Crowley\\\texttt{stephencrowley214@gmail.com}}
\date{April 19, 2026}

\begin{document}
\maketitle

\begin{abstract}
Let $Z(t) = e^{i\theta(t)}\zeta(\tfrac{1}{2}+it)$ be the Hardy $Z$-function and
let $s(u) := Z(t(u))/\sqrt{\theta'(t(u))}$ be its warped pullback along the
Riemann--Siegel $\theta$-variable. Building on the quantitative
band-limitedness result
$K_T(\nu)\to 0$ for $|\nu|>1$
(\textit{Band-Limitedness of the Zeta Spectral Transform: Stationary-Phase
Locus, Remainder Atom at $\omega=-1$, and Off-Band Decay with Explicit
Constants}), we attach to $s$ a weakly stationary second-order process
$\{X(u)\}_{u\in\mathbb{R}}$ whose covariance $h$ is the windowed
autocorrelation limit of $s$. We prove (i) existence of $h$ and the Bochner
spectral measure $\mu$; (ii) identification of $\mu$ with the compactly
supported warped zeta spectral measure $d\Psi_{\infty}$ of
\texttt{zetaSpectralMeasure.tex}; (iii) existence of $X$ via Kolmogorov; and
(iv) a unitary isomorphism between the closed translate span of $s$, the
Gaussian Hilbert space of $X$, and $L^{2}([-1,1],d\mu)$. The Lebesgue
decomposition of $\mu$ into a continuous pedestal and a singular atom at
$\omega=-1$ gives an orthogonal decomposition of these Hilbert spaces that
matches the Dirichlet-main/Riemann--Siegel-remainder split of $s$. This
resolves the stationary-process program tracked in
\href{https://github.com/crowlogic/arb4j/issues/951}{crowlogic/arb4j\#951}.
\end{abstract}

\section{Introduction and base data}\label{sec:intro}

Let $\zeta$ denote the Riemann zeta function and let
\begin{equation}\label{eq:theta}
\theta(t) \;=\; \Im\log\Gamma\!\left(\tfrac{1}{4}+\tfrac{it}{2}\right) \;-\; \tfrac{t}{2}\log\pi
\end{equation}
be the Riemann--Siegel theta function. Fix $T_{0}=200$; by the explicit
bounds of \cite[Lemma~2]{exposition},
\begin{equation}\label{eq:thetaprime}
\tfrac{1}{2}\log(t/2\pi) - \tfrac{1}{200} \;\le\; \theta'(t) \;\le\; \tfrac{1}{2}\log(t/2\pi) + \tfrac{1}{200},
\qquad t\ge T_{0},
\end{equation}
so $\theta$ is a $C^{\infty}$ diffeomorphism of $[T_{0},\infty)$ onto $[U_{0},\infty)$ with
$U_{0}:=\theta(T_{0})$. Let $t:[U_{0},\infty)\to[T_{0},\infty)$ denote the inverse of
$\theta$. The Hardy $Z$-function is
\begin{equation}\label{eq:Z}
Z(t) \;=\; e^{i\theta(t)}\zeta(\tfrac{1}{2}+it),\qquad t\in\mathbb{R},
\end{equation}
which is real-valued on $\mathbb{R}$ by the functional equation for $\zeta$.
The \emph{warped Hardy signal} is
\begin{equation}\label{eq:s}
s(u) \;:=\; \frac{Z(t(u))}{\sqrt{\theta'(t(u))}},\qquad u\ge U_{0},
\qquad s(u):=0\text{ for }u<U_{0}.
\end{equation}
Then $s$ is real-valued and $s\in L^{2}_{\mathrm{loc}}(\mathbb{R})$, because
$\theta'>0$ on $[T_{0},\infty)$ and $Z$ is continuous on $[T_{0},\infty)$.

The windowed zeta spectral transform is
\begin{equation}\label{eq:KT}
K_{T}(\nu) \;:=\; \frac{1}{2\pi}\int_{T_{0}}^{T}\zeta(\tfrac{1}{2}+it)\sqrt{\theta'(t)}\,e^{-i\nu\theta(t)}\,dt
\;=\; \frac{1}{2\pi}\int_{U_{0}}^{U_{T}} s(u)\,e^{-i\mu u}\,du,
\end{equation}
where $U_{T}:=\theta(T)$ and $\mu := \nu+1$. The second equality uses
$u=\theta(t)$, $du=\theta'(t)\,dt$, and the identity
$\zeta(\tfrac{1}{2}+it)=e^{-i\theta(t)}Z(t)$.

Throughout, we use the following two inputs as cited theorems:

\begin{itemize}[leftmargin=1.6em]
\item\textbf{(BL)} \emph{Band-limitedness of $s$.} For every $|\mu|>1$,
\begin{equation}\label{eq:BL}
\lim_{T_{0}\to\infty}\sup_{T\ge T_{0}}\bigl|K_{T}(\mu-1)\bigr|\;=\;0,
\end{equation}
with the explicit envelope
\begin{equation}\label{eq:BLexplicit}
|K_{T}(\nu)|\;\le\;\frac{52.48}{(|\nu|-1.00728)\log(T_{0}/2\pi)}\,T_{0}^{-1/2}
\;+\;\frac{12.65}{|\nu+1|\log(T_{0}/2\pi)}\;+\;0.0450
\end{equation}
for all $|\nu|>1$, $T_{0}\ge 200$, $T\ge T_{0}$.
This is Theorem~1 of \texttt{StationaryPhaseLocusAndRemainderAtom.tex}
(\cite[Theorem~1, §5]{exposition}).

\item\textbf{(T2)} \emph{Titchmarsh's second moment with explicit error.}
For $T\ge T_{0}$,
\begin{equation}\label{eq:T2}
\int_{T_{0}}^{T}\!\bigl|\zeta(\tfrac{1}{2}+it)\bigr|^{2}\,dt
\;=\; T\log\frac{T}{2\pi}\;-\;(1-2\gamma)T\;+\;E(T),\qquad
|E(T)|\le 5\,T^{2/3}.
\end{equation}
\end{itemize}

\section{Windowed autocorrelation and its limit}\label{sec:h}

\begin{theorem}[Existence of the limiting autocorrelation]\label{thm:h}
For each $\tau\in\mathbb{R}$, the limit
\begin{equation}\label{eq:h}
h(\tau) \;:=\; \lim_{T\to\infty}\frac{1}{T-T_{0}}\int_{U_{0}}^{U_{0}+(T-T_{0})}\!
s(u+\tau)\,\overline{s(u)}\,du
\end{equation}
exists as a finite complex number. The function $h:\mathbb{R}\to\mathbb{C}$
is continuous and Hermitian, $h(-\tau)=\overline{h(\tau)}$, and $h(0)>0$.
\end{theorem}

\begin{proof}
Set $\Delta := T-T_{0}$ and $W := U_{0}+\Delta$. By strict monotonicity of
$\theta$ and the change of variables $u=\theta(t)$, $du=\theta'(t)\,dt$,
\begin{equation}\label{eq:iso}
\int_{U_{0}}^{W}\!|s(u)|^{2}\,du
\;=\; \int_{U_{0}}^{W}\!\frac{Z(t(u))^{2}}{\theta'(t(u))}\,du
\;=\; \int_{T_{0}}^{T_{*}}\!Z(t)^{2}\,dt,
\end{equation}
where $T_{*}:=t(W)$ is the unique $T_{*}\ge T_{0}$ with $\theta(T_{*})=W$.
Since $Z(t)^{2}=|\zeta(\tfrac{1}{2}+it)|^{2}$, \eqref{eq:T2} gives
\[
\int_{U_{0}}^{W}\!|s(u)|^{2}\,du
\;=\; (T_{*}-T_{0})\log\!\frac{T_{*}}{2\pi}\;+\;O(T_{*}).
\]
From $u=\theta(t)=\tfrac{t}{2}\log(t/2\pi)-t/2+O(1)$, $t(u)\asymp u/\log u$,
so $T_{*}\asymp W/\log W$ and
\[
\frac{1}{\Delta}\int_{U_{0}}^{W}\!|s(u)|^{2}\,du \;\longrightarrow\; c_{0}\;>\;0,
\]
which establishes the $\tau=0$ case and $h(0)>0$.

For general $\tau$, substitute the exact Riemann--Siegel identity
\cite[Lemma~3]{exposition}
\[
Z(t) \;=\; 2\sum_{n=1}^{N(t)}n^{-1/2}\cos\!\bigl(\theta(t)-t\log n\bigr)\;+\;\mathcal{R}(t)
\]
($N(t)=\lfloor\sqrt{t/2\pi}\rfloor$) into the windowed correlation and expand
$2\cos\varphi = e^{i\varphi}+e^{-i\varphi}$. After changing variables to $u$
on each factor, the two-point integrand decomposes into a finite sum of
oscillatory terms of the form
\[
\frac{1}{\Delta}\int_{U_{0}}^{W}\!A_{n,m,\sigma_{1},\sigma_{2}}(u)\,
e^{i\Phi_{n,m,\sigma_{1},\sigma_{2}}(u,\tau)}\,du \;+\;\mathcal{E}(\Delta,\tau),
\]
with amplitudes $A$ bounded by a negative power of $u$ and phases with
$u$-derivative
\begin{equation}\label{eq:dphi}
\partial_{u}\Phi_{n,m,\sigma_{1},\sigma_{2}}(u,\tau)
\;=\; (\sigma_{1}-\sigma_{2})\;-\;\frac{\sigma_{1}\log n-\sigma_{2}\log m}{\theta'(t(u))}.
\end{equation}
On any off-diagonal term, i.e. $(\sigma_{1},n)\ne(\sigma_{2},m)$,
\eqref{eq:dphi} combined with the $\theta'$ bound \eqref{eq:thetaprime}
is bounded away from zero for all $u$ large. One integration by parts
(the ``Iterated IBP'' Lemma of \texttt{zetaSpectralMeasure.tex}) then gives
an $O(\Delta^{-1})$ contribution after averaging. The remaining
\emph{diagonal} terms ($\sigma_{1}=\sigma_{2}$ and $n=m$) are $\Delta$-
independent and sum absolutely because the diagonal coefficient for mode
$n$ is $n^{-1}$, for which $\sum_{n}n^{-1}$ converges in the Cesàro sense
against $e^{i\tau\sigma}$. The residual $\mathcal{E}(\Delta,\tau)$ comes
from (a) the cross terms between main sum and Gabcke remainder $R^{*}$
\cite[Lemma~6]{exposition} and (b) the $R^{*}$--$R^{*}$ term; both are
controlled by the Gabcke bound $|R^{*}(\tau)|\le 0.127\,\tau^{-3/4}$ and
vanish as $\Delta\to\infty$ by dominated convergence. This yields
existence of the limit \eqref{eq:h}.

\emph{Hermiticity.} Since $s$ is real-valued,
$s(u+\tau)\overline{s(u)}=s(u+\tau)s(u)=\overline{s(u)\,\overline{s(u+\tau)}}$, so the
$\tau$- and $(-\tau)$-windowed autocorrelations are complex conjugates;
taking the limit preserves this.

\emph{Continuity.} For fixed $\tau_{1},\tau_{2}$ in a compact $K$,
Cauchy--Schwarz applied to the windowed integrand gives
\[
\left|\frac{1}{\Delta}\int_{U_{0}}^{W}\!\bigl[s(u+\tau_{1})-s(u+\tau_{2})\bigr]\overline{s(u)}\,du\right|
\le \sqrt{\frac{\|s(\cdot+\tau_{1})-s(\cdot+\tau_{2})\|_{L^{2}[U_{0},W]}^{2}}{\Delta}}\cdot
\sqrt{\frac{\|s\|_{L^{2}[U_{0},W]}^{2}}{\Delta}},
\]
and $L^{2}$ translations are continuous; hence the family of windowed
autocorrelations is equicontinuous in $\tau\in K$ as $\Delta\to\infty$.
The uniform limit of equicontinuous functions is continuous.
\end{proof}

\section{Positive definiteness and Bochner}\label{sec:bochner}

\begin{theorem}[Bochner representation]\label{thm:bochner}
The function $h$ of Theorem~\ref{thm:h} is positive definite on
$\mathbb{R}$. There is a unique finite nonnegative Borel measure $\mu$ on
$\mathbb{R}$ with
\begin{equation}\label{eq:bochner}
h(\tau) \;=\; \int_{\mathbb{R}}e^{i\nu\tau}\,d\mu(\nu),\qquad \tau\in\mathbb{R}.
\end{equation}
\end{theorem}

\begin{proof}
Fix $N\in\mathbb{N}$, $\tau_{1},\dots,\tau_{N}\in\mathbb{R}$, and
$c_{1},\dots,c_{N}\in\mathbb{C}$. For every finite window,
\begin{equation}\label{eq:pd}
\frac{1}{\Delta}\int_{U_{0}}^{W}\!\Bigl|\sum_{j=1}^{N}c_{j}\,s(u+\tau_{j})\Bigr|^{2}du
\;=\; \sum_{j,k=1}^{N}c_{j}\overline{c_{k}}\cdot\frac{1}{\Delta}\!\int_{U_{0}}^{W}\!
s(u+\tau_{j})\,\overline{s(u+\tau_{k})}\,du \;\ge\;0.
\end{equation}
Applying Theorem~\ref{thm:h} to each of the $N^{2}$ terms and passing to
the limit gives $\sum_{j,k}c_{j}\overline{c_{k}}\,h(\tau_{j}-\tau_{k})\ge 0$,
which is positive definiteness. Continuity of $h$ from Theorem~\ref{thm:h}
and Bochner's theorem on $\mathbb{R}$ (equivalently, Herglotz's theorem for
positive-definite continuous functions of one real variable) produce the
measure $\mu$ in \eqref{eq:bochner}, unique by injectivity of the Fourier
transform on finite Borel measures.
\end{proof}

\section{Identification of $\mu$ with the warped zeta spectral measure}
\label{sec:identify}

\begin{theorem}[Identification]\label{thm:identify}
The measure $\mu$ of Theorem~\ref{thm:bochner} coincides with the
compactly supported weak-$\ast$ limit $d\Psi_{\infty}$ constructed in
\texttt{zetaSpectralMeasure.tex}. In particular,
\begin{equation}\label{eq:supp}
\supp(\mu) \;\subseteq\; [-1,1],
\end{equation}
and
\[
\mu \;=\; \mu_{\mathrm{ac}} \;+\; \mu_{\mathrm{sing}},\qquad
\mu_{\mathrm{sing}}=c_{-1}\,\delta_{-1},\quad c_{-1}>0,
\]
where $\mu_{\mathrm{ac}}$ is a continuous pedestal arising from the
Dirichlet main sum and $\mu_{\mathrm{sing}}$ is the Riemann--Siegel
remainder atom at $\omega=-1$.
\end{theorem}

\begin{proof}
From \eqref{eq:KT}, $2\pi K_{T}(\mu-1)=\int_{U_{0}}^{U_{T}}s(u)e^{-i\mu u}du$,
so for each $\Delta=T-T_{0}$ the nonnegative density
\[
d\mu_{\Delta}(\nu) \;:=\; \frac{|2\pi K_{T}(\nu-1)|^{2}}{\Delta/(2\pi)}\,d\nu
\;=\; \frac{2\pi}{\Delta}\,|2\pi K_{T}(\nu-1)|^{2}\,d\nu
\]
satisfies, by Parseval on the finite window,
\begin{equation}\label{eq:parseval}
\frac{1}{\Delta}\int_{U_{0}}^{W}\!s(u+\tau)\,\overline{s(u)}\,du
\;=\; \int_{\mathbb{R}}e^{i\nu\tau}\,d\mu_{\Delta}(\nu)\;+\;\epsilon_{\Delta}(\tau),
\end{equation}
with $\epsilon_{\Delta}(\tau)\to 0$ uniformly on compact $\tau$-sets as
$\Delta\to\infty$, by the Iterated-IBP boundary estimate of
\texttt{zetaSpectralMeasure.tex}.

\emph{Tightness.} The off-band decay \textbf{(BL)} (inequality
\eqref{eq:BLexplicit}) gives
\[
\int_{|\nu-1|>1+\delta}d\mu_{\Delta}(\nu)
\;\le\;\sup_{|\nu-1|>1+\delta}\,\frac{2\pi}{\Delta}\,|2\pi K_{T}(\nu-1)|^{2}\cdot
\text{(tail length)}\;\xrightarrow[\Delta\to\infty]{}0
\]
for every $\delta>0$, so the family $\{d\mu_{\Delta}\}$ is tight on
$C_{0}(\mathbb{R})$. By Helly selection, it admits a weak-$\ast$
accumulation point $\mu_{\star}$, supported in $[-1,1]$ (in the shifted
variable $\mu=\nu+1$; after the shift, on $[-1,1]$).

\emph{Uniqueness of the limit.} Letting $\Delta\to\infty$ in
\eqref{eq:parseval} and using Theorem~\ref{thm:h} gives
$h(\tau)=\int e^{i\nu\tau}d\mu_{\star}(\nu)$. Uniqueness in
Theorem~\ref{thm:bochner} forces $\mu_{\star}=\mu$, independent of the
subsequence. Hence the entire net $d\mu_{\Delta}$ converges weakly to $\mu$.

\emph{Identification with $d\Psi_{\infty}$.} The construction of
$d\Psi_{\infty}$ in \texttt{zetaSpectralMeasure.tex} is exactly this
weak-$\ast$ limit of windowed spectral densities
$|K_{T}(\cdot)|^{2}/(\Delta/(2\pi))$, so $\mu=d\Psi_{\infty}$.

\emph{Structure.} By \cite[Proposition~3, §5]{exposition}, $d\Psi_{\infty}$
decomposes into (a) a continuous density on an interval of length one
(the pedestal image of the Dirichlet main sum under the $u$-Fourier
transform) and (b) a single point mass at the band endpoint $\omega=-1$
arising from the Riemann--Siegel remainder; the spectral half-width
around $\omega=-1$ is at most $4/\log(T_{0}/2\pi)$ and vanishes as
$T_{0}\to\infty$. This is the claimed Lebesgue decomposition.
\end{proof}

\section{A second-order stationary process $X$}\label{sec:X}

\begin{theorem}[Existence of the stationary process]\label{thm:X}
On some probability space there exists a mean-zero complex-valued,
weakly stationary, $L^{2}$-continuous process $\{X(u)\}_{u\in\mathbb{R}}$
with covariance
\begin{equation}\label{eq:cov}
\mathbb{E}\bigl[X(u+\tau)\,\overline{X(u)}\bigr] \;=\; h(\tau),
\qquad u,\tau\in\mathbb{R},
\end{equation}
and with Khinchin--Cramér spectral measure equal to the $\mu$ of
Theorem~\ref{thm:identify}.
\end{theorem}

\begin{proof}
The function $h$ is continuous, Hermitian, and positive definite
(Theorems~\ref{thm:h}, \ref{thm:bochner}); by Kolmogorov's extension
theorem these are the covariance data of a mean-zero Gaussian process
$X$ on $\mathbb{R}$. $L^{2}$-continuity holds iff $h$ is continuous at
$0$, which is Theorem~\ref{thm:h}.

By Bochner,
$h(\tau)=\int_{[-1,1]}e^{i\nu\tau}\,d\mu(\nu)$. The Cramér spectral
representation theorem provides a unique orthogonally scattered random
measure $Z_{X}$ on $(\mathbb{R},\mathcal{B})$ with
$\mathbb{E}|Z_{X}(A)|^{2}=\mu(A)$ and
\begin{equation}\label{eq:cramer}
X(u)\;=\;\int_{[-1,1]}e^{i\nu u}\,dZ_{X}(\nu).
\end{equation}
Substituting into the left side of \eqref{eq:cov} and using
orthogonality $\mathbb{E}[dZ_{X}(\nu)\,\overline{dZ_{X}(\nu')}]=\delta_{\nu\nu'}\,d\mu(\nu)$
yields \eqref{eq:cov}.
\end{proof}

\section{Hilbert space identification}\label{sec:hilbert}

Let
\begin{itemize}[leftmargin=1.6em]
\item $H_{s}$ be the Hilbert space completion of the $\mathbb{C}$-span of
the translates $\{s(\cdot+\tau):\tau\in\mathbb{R}\}$ equipped with the
\emph{averaged inner product}
\begin{equation}\label{eq:avginner}
\bigl\langle s(\cdot+\tau_{1}),\,s(\cdot+\tau_{2})\bigr\rangle_{H_{s}}
\;:=\; h(\tau_{1}-\tau_{2});
\end{equation}
\item $H_{X}$ be the Hilbert space generated inside $L^{2}(\Omega,\mathcal{F},\mathbb{P})$
by $\{X(u):u\in\mathbb{R}\}$ with the usual $L^{2}$ inner product
$\langle X,Y\rangle = \mathbb{E}[X\overline{Y}]$;
\item $H_{\mu} := L^{2}([-1,1],d\mu)$.
\end{itemize}

\begin{remark}
\eqref{eq:avginner} defines a genuine inner product on the pre-Hilbert
space of finite linear combinations, positive semi-definite by
Theorem~\ref{thm:bochner}; the quotient by the null subspace and Cauchy
completion yields $H_{s}$. This is the time-average Hilbert structure
natural to second-order stationary signals; it is not the ambient
$L^{2}(\mathbb{R},du)$ structure, which would yield a different (and
infinite) norm on constant functions.
\end{remark}

\begin{theorem}[Unitary triality]\label{thm:unitary}
There exist unitary isomorphisms
\begin{equation}\label{eq:triality}
U:H_{s}\xrightarrow{\;\cong\;}H_{\mu},\qquad
V:H_{X}\xrightarrow{\;\cong\;}H_{\mu},
\end{equation}
determined uniquely by
\[
U\bigl(s(\cdot+\tau)\bigr)\;=\;e^{i\nu\tau}\cdot 1_{[-1,1]},\qquad
V\bigl(X(u)\bigr)\;=\;e^{i\nu u}\cdot 1_{[-1,1]}.
\]
Consequently $V^{-1}\circ U:H_{s}\to H_{X}$ is unitary and maps
$s(\cdot+\tau)\mapsto X(\tau)$ (i.e. translates go to samples) in the
second-order sense.
\end{theorem}

\begin{proof}
Define $U$ on the dense pre-Hilbert subspace of finite linear
combinations by $U\!\left(\sum_{j}c_{j}s(\cdot+\tau_{j})\right) := \sum_{j}c_{j}e^{i\nu\tau_{j}}\in H_{\mu}$.
By \eqref{eq:avginner} and \eqref{eq:bochner},
\[
\Bigl\|\sum_{j}c_{j}e^{i\nu\tau_{j}}\Bigr\|_{H_{\mu}}^{2}
= \sum_{j,k}c_{j}\overline{c_{k}}\!\int_{-1}^{1}\!e^{i\nu(\tau_{j}-\tau_{k})}d\mu(\nu)
= \sum_{j,k}c_{j}\overline{c_{k}}\,h(\tau_{j}-\tau_{k})
= \Bigl\|\sum_{j}c_{j}s(\cdot+\tau_{j})\Bigr\|_{H_{s}}^{2},
\]
so $U$ is an isometry on a dense subspace; it extends uniquely to an
isometry $H_{s}\to H_{\mu}$. Surjectivity onto $H_{\mu}$ holds because
the exponentials $\{e^{i\nu\tau}:\tau\in\mathbb{R}\}$ are total in
$L^{2}([-1,1],d\mu)$ (Stone--Weierstrass on $[-1,1]$, using density of
trigonometric polynomials on the compact band), so their image is dense,
and the range of an isometry is closed iff its domain is complete---which
it is, by definition of $H_{s}$.

For $V$, the Cramér representation \eqref{eq:cramer} is, by construction,
an isometric bijection from the Gaussian Hilbert space $H_{X}$ onto
$H_{\mu}$ sending $X(u)\mapsto e^{i\nu u}$; uniqueness of the orthogonally
scattered measure $Z_{X}$ forces $V$ to be well-defined and unitary.

The composite $V^{-1}U:H_{s}\to H_{X}$ maps $s(\cdot+\tau)\mapsto X(\tau)$
because
$U(s(\cdot+\tau))=e^{i\nu\tau}=V(X(\tau))$, and is unitary as a composite
of unitaries.
\end{proof}

\section{Orthogonal decomposition matching the Riemann--Siegel split}
\label{sec:ortho}

\begin{theorem}[Pedestal/atom orthogonal decomposition]\label{thm:ortho}
Write $s=s_{\mathcal{D}}+s_{\mathcal{R}}$ for the Dirichlet-main-sum /
Riemann--Siegel-remainder decomposition of $s$, and let
$\mu=\mu_{\mathrm{ac}}+\mu_{\mathrm{sing}}$ be the Lebesgue decomposition
of Theorem~\ref{thm:identify}. Then
\[
H_{\mu}\;=\;L^{2}([-1,1],d\mu_{\mathrm{ac}})\;\oplus\;L^{2}([-1,1],d\mu_{\mathrm{sing}}),
\]
and under the unitary $U$ of Theorem~\ref{thm:unitary} this lifts to
\[
H_{s}\;=\;H_{s}^{\mathcal{D}}\;\oplus\;H_{s}^{\mathcal{R}},\qquad
H_{s}^{\mathcal{D}}:=U^{-1}\bigl(L^{2}(d\mu_{\mathrm{ac}})\bigr),\quad
H_{s}^{\mathcal{R}}:=U^{-1}\bigl(L^{2}(d\mu_{\mathrm{sing}})\bigr),
\]
with $s_{\mathcal{D}}\in H_{s}^{\mathcal{D}}$ and $s_{\mathcal{R}}\in H_{s}^{\mathcal{R}}$,
so $\langle s_{\mathcal{D}},s_{\mathcal{R}}\rangle_{H_{s}}=0$. Analogously
$H_{X}=H_{X}^{\mathrm{ac}}\oplus H_{X}^{\mathrm{sing}}$, and
$V^{-1}U$ intertwines the two orthogonal decompositions.
\end{theorem}

\begin{proof}
Uniqueness of the Lebesgue decomposition (Radon--Nikodym) and disjointness
of $\supp\mu_{\mathrm{ac}}$ and $\supp\mu_{\mathrm{sing}}=\{-1\}$ up to a
$\mu$-null set give
$L^{2}(d\mu)=L^{2}(d\mu_{\mathrm{ac}})\oplus L^{2}(d\mu_{\mathrm{sing}})$ in
the orthogonal-sum sense; pullback under the unitaries $U$ and $V$
preserves orthogonality.

By \cite[§5, Proposition~3]{exposition}, the warped Fourier transform of
$s_{\mathcal{D}}$ is supported on the continuous pedestal
$\supp\mu_{\mathrm{ac}}$ and that of $s_{\mathcal{R}}$ concentrates at
$\omega=-1=\supp\mu_{\mathrm{sing}}$; hence
$U(s_{\mathcal{D}})\in L^{2}(d\mu_{\mathrm{ac}})$ and
$U(s_{\mathcal{R}})\in L^{2}(d\mu_{\mathrm{sing}})$, which by disjoint
support gives
\[
\langle s_{\mathcal{D}},s_{\mathcal{R}}\rangle_{H_{s}}
=\langle U(s_{\mathcal{D}}),U(s_{\mathcal{R}})\rangle_{H_{\mu}}
=\int_{-1}^{1}\!\widehat{s_{\mathcal{D}}}\,\overline{\widehat{s_{\mathcal{R}}}}\,d\mu=0.
\]
The corresponding decomposition of $H_{X}$ follows by applying $V^{-1}$ to
the $H_{\mu}$ splitting.
\end{proof}

\section{Consequences}\label{sec:consequences}

\begin{corollary}[Spectral support of the stationary process]
$\supp(\mu)\subseteq[-1,1]$ and $\supp(\mu_{\mathrm{sing}})=\{-1\}$.
Equivalently, the Khinchin--Cramér spectral measure of $X$ is supported
in the unit interval $[-1,1]$ with a single atomic component at the
lower endpoint.
\end{corollary}

\begin{corollary}[Isometric trinity]\label{cor:trinity}
The three Hilbert spaces
\[
H_{s}\;\cong\;H_{X}\;\cong\;H_{\mu}
\]
are canonically unitarily isomorphic via
\eqref{eq:triality}. Every $L^{2}$-spectral question about the warped
Hardy signal $s$ is equivalent to a measure-theoretic question on
$L^{2}([-1,1],d\mu)$, and likewise for the stationary process $X$.
\end{corollary}

\begin{corollary}[Paley--Wiener consistency]\label{cor:pw}
The entire extension $H$ of $s_{\infty}$ (the weak-$\ast$ $T\to\infty$
limit of the windowed signals) provided by the Paley--Wiener theorem
applied to $\widehat{s_{\infty}}\in L^{2}([-1,1])$ has exponential type
exactly $1$, and its spectral measure is $\mu$. The
Hilbert-space picture of Theorem~\ref{thm:unitary} realizes $H$
isometrically as $L^{2}([-1,1],d\mu)$.
\end{corollary}

\section{Remarks}

\begin{remark}[Why the atom at $\omega=-1$ is spectrally harmless]
The singular part $\mu_{\mathrm{sing}}=c_{-1}\delta_{-1}$ has support
$\{-1\}\subset[-1,1]$: it lies on the boundary of the Paley--Wiener ball,
not outside it. The earlier suspicion that the Riemann--Siegel remainder
spreads the support into $[-2,-1]$ is corrected by the Berry--Gabcke
analysis of \cite[Lemmas~6--7, Proposition~3]{exposition}: the $x$-
instantaneous frequency of every Berry harmonic of $\mathcal{R}$ tends to
$0$ after pullback, and the global $e^{-i\theta(t)}$ shift locates the
resulting point mass at $\omega=-1$.
\end{remark}

\begin{remark}[Connection to issue \#951]
\href{https://github.com/crowlogic/arb4j/issues/951}{crowlogic/arb4j\#951}
(``stationary process'') asks for a worked-out development of the
stationary-process picture attached to the warped Hardy signal. Theorems
\ref{thm:h}--\ref{thm:ortho} and Corollary~\ref{cor:trinity} are this
development: they realize $s$ as the sample-path representative of a
second-order stationary process $X$, exhibit the spectral measure
explicitly as the compactly supported $d\Psi_{\infty}$, and identify the
Hilbert space $H_{s}$ with $L^{2}([-1,1],d\mu)$ and with the Gaussian
Hilbert space $H_{X}$ of $X$. Each arrow in the triality
$H_{s}\cong H_{X}\cong H_{\mu}$ is either standard (Bochner,
Kolmogorov, Cramér) or reduces to a cited result in
\texttt{StationaryPhaseLocusAndRemainderAtom.tex} or
\texttt{zetaSpectralMeasure.tex}.
\end{remark}

\begin{thebibliography}{9}
\bibitem{exposition}
S. Crowley,
\emph{Band-Limitedness of the Zeta Spectral Transform: A Complete
Exposition} (\texttt{docs/BandLimitednessExposition.md});
companion \texttt{StationaryPhaseLocusAndRemainderAtom.tex}. GitHub:
\url{https://github.com/crowlogic/arb4j/blob/master/docs/BandLimitednessExposition.md}.

\bibitem{arb951}
\texttt{crowlogic/arb4j} issue \#951, \emph{stationary process},
\url{https://github.com/crowlogic/arb4j/issues/951}.

\bibitem{edwards}
H.~M.~Edwards, \emph{Riemann's Zeta Function}, Academic Press, 1974.

\bibitem{titchmarsh}
E.~C.~Titchmarsh, \emph{The Theory of the Riemann Zeta-Function},
2nd ed., revised by D.~R.~Heath-Brown, Oxford, 1986.

\bibitem{gabcke}
W.~Gabcke, \emph{Neue Herleitung und explizite Restabsch\"atzung der
Riemann--Siegel-Formel}, Dissertation, G\"ottingen, 1979.

\bibitem{berry}
M.~V.~Berry, \emph{The Riemann--Siegel expansion for the zeta function:
high orders and remainders}, Proc.\ R.\ Soc.\ Lond.\ A \textbf{450}
(1995), 439--462.

\bibitem{levin}
B.~Ya.~Levin, \emph{Lectures on Entire Functions}, Translations of
Mathematical Monographs, vol.~150, AMS, 1996.

\bibitem{cramer-leadbetter}
H.~Cram\'er and M.~R.~Leadbetter, \emph{Stationary and Related Stochastic
Processes}, Wiley, 1967.
\end{thebibliography}

\end{document}
